\section{Intersections of divisor-degree images}
\label{sec:degree-intersections-v152}

Here the total relation degree $D$ is fixed.  Write $Y_a\subset
\Gr(r,S_D)$ for the image admitting a divisor of degree $a$;
empty Grassmannians are omitted.  In several variables these images
are not generally nested.  Their reduced intersections nevertheless
have a finite description by common parts of two marked divisors.

\begin{proposition}[The common-part formula]
\label{prop:degree-intersections-v152}
For $1\leq a,b<D$ and $0\leq j\leq\min(a,b)$, let $M_j$ be the
scheme-theoretic image of
\begin{equation}\label{eq:intersection-map-v152}
\begin{split}
 \Pj(S_j)\times\Pj(S_{a-j})\times\Pj(S_{b-j})
       \times\Gr(r,S_{D-a-b+j})&\longrightarrow\Gr(r,S_D),\\
                   (c,u,v,L)&\longmapsto cuvL.
\end{split}
\end{equation}
A negative last degree or insufficient dimension gives the empty
image, and $\Pj(S_0)$ is a point.  Then
\begin{equation}\label{eq:intersection-reduced-v152}
             (Y_a\cap Y_b)_{\mathrm{red}}
                         =\left(\bigcup_j M_j\right)_{\mathrm{red}}.
\end{equation}
On the locus of two marked divisors $f_a,f_b$ with exact common
part of degree $j$, the corresponding factorization is
$c=\gcd(f_a,f_b)$, $f_a=cu$, $f_b=cv$, with $\gcd(u,v)=1$.
In particular, at a relation $K=a_0^mL$ with $\gcd(L)=1$, the
set of positive divisor degrees is exactly $\{1,\ldots,m\}$.
For the boundary point in Theorem~\ref{thm:universal-boundary-v152},
$m=D-q_n$.
\end{proposition}
\begin{proof}
Each map in \eqref{eq:intersection-map-v152} is defined on a
projective scheme: multiplication by a nonzero form has constant
rank, so the subspace construction is a morphism.  Its image is
closed and admits the two divisors $cu$ and $cv$, proving one
inclusion.  Conversely, if $K$ has divisors $f_a,f_b$, unique
factorization gives their common part $c$ of some degree $j$.
Their least common multiple is $cuv$ and divides every member of
$K$.  Division yields $L$ of degree $D-a-b+j$, of the same rank.
This proves equality on geometric points.  Over $\C$, a reduced
closed subscheme of a finite-type scheme is determined by its closed
points, giving \eqref{eq:intersection-reduced-v152}.  The last
assertion follows because every homogeneous divisor of a power
of one linear form is a power of that form.
\end{proof}

The reduction in \eqref{eq:intersection-reduced-v152} is essential
to the statement of this global formula.  It is not a prescription
to reduce the individual $Y_a$, their normalization fibres, or their
singular subschemes.  The full intersection of the two formal
branches on the coprime stratum was computed over Artin bases in
Theorem~\ref{thm:split-conductor-v152};
Theorem~\ref{thm:binary-pinch-v152} computes the thickening when
those branches collide.  Thus the degree-poset description and
the local scheme calculations have distinct, explicit roles.
